\section{闭合而不封口：接收一个区域世界}
\label{sec:jin-unification}

二六三年至二八〇年的胜利把蜀、魏、吴的最高权属接入西晋，却没有让三处中枢
的旧问题消失。灭蜀后的\eventanchor{TK-0264-ZHONG-HUI-REBELLION}{成都兵变}首先
提醒征服者：降表并不等于掌握降军、道路和地方信息。魏晋禅代说明近畿军政
控制可以先于名义交接；灭吴则说明长江防线必须被多路水陆协同、荆州经营和
益州船队同时穿透。

\subsection{三种接收速度}

\textbf{诏册的速度}最快。受禅与降表可以在很短时间内改变最高名义。
\textbf{军政节点的速度}较慢：成都、襄阳、西陵和建业都需要新的命令链、守备
和地方协作。\textbf{地方生活的速度}最慢：县廷的簿籍、田租、役使、港口和
家族关系不会因为一次典礼而自动换成同一种实践。把三种速度分开，才能理解
为什么统一既是事实，又是一段漫长的工作。

二七六年杜预受命参与荆州军务，
\eventanchor{TK-DENS-0276-01}{西晋增加了一名熟悉军政的伐吴统帅}；这只说明
统帅人选和命令链在调整，不能由任命反推完整作战方案。二八〇年孙皓出降后被
送往洛阳，西晋以封爵安置末帝与部分降臣，
\eventanchor{TK-DENS-0280-01}{接收也有仪式与安置的一面}。帝纪能够确认降服
与封爵的顺序，无法据此证明所有吴地官员、港口和家族已经被同样地接入新秩序。

\subsection{本卷的收束}

共同时间线从曹操死后洛阳的接收开始，到孙皓出降结束；三个 dossier 显示，
魏的近畿中枢、蜀汉的山道后方和孙吴的江面网络各有不同的资源结构；跨政权
综合则说明道路、盟约和边缘持续改写这些结构。西晋的统一不是跳出这张网络，
而是以更大的尺度继承它，并把新的制度承诺送往仍然不等速的地方。

\begin{debate}[“统一”的边界]
《晋书》帝纪与《资治通鉴》都给二八〇年以明确的政治终点。这个终点可用于
叙述全国最高政权的合并；若把它扩展为行政、经济或社会已经无摩擦地一体化，
便超出了这些材料所能证明的范围。
\end{debate}


\subsection{晋纪补链：统一之前的诏令、接收与边防}

以下晋纪补链把统一前的诏令、任免、边防与接收节点接到三种不同速度的收束之中。它们确认的是本纪所载的命令或结果，不能据此预定地方执行、降人选择或统一后的社会整合。

\paragraph{景元四年（263年）} 《本纪》在景元四年（263年）先把这些动作排入同一条年序：\eventanchor{TK-0263-WEI-INVADES-SHU}{司马昭命钟会、邓艾、诸葛绪多路伐蜀，魏军由关中和陇右南下}。\footnote{核心来源：《三国志·邓艾传》；《三国志·钟会传》；《三国志·姜维传》；《三国志·司马昭传》；《资治通鉴》魏纪。材料限度：出师、主要统帅和多路部署可靠；兵数、诸葛绪失职、阴平路线筹划及各军进度有争议}

\paragraph{咸熙二年八月（265年）} 到了咸熙二年八月（265年），朝廷可见的变化并不只在一项大事上，\eventanchor{TK-0265-SIMA-ZHAO-DEATH}{司马昭去世，司马炎继承晋王位及魏廷实际军政控制权}。\footnote{核心来源：《三国志·陈留王纪》；《晋书·文帝纪》；《晋书·武帝纪》。材料限度：卒年与继承可靠；遗命、司马炎接掌禁军和魏帝朝臣的协商空间有争议}

\paragraph{泰始五年（269年）} 记录写下泰始五年（269年）的多个接口：\eventanchor{TK-0269-YANG-HU-JINGZHOU}{羊祜都督荆州诸军事，在襄阳一线经营屯戍、招纳与对吴关系}。\footnote{核心来源：《晋书·羊祜传》；《三国志·陆抗传》；《资治通鉴》晋纪。材料限度：任命与荆州经营可靠；招降规模、怀柔政策效果和羊祜个人作用有争议}

\paragraph{咸宁四年（278年）} 咸宁四年（278年）的条目把命令、任免或接触并列起来：\eventanchor{TK-0278-YANG-HU-DEATH}{羊祜去世前上表建议伐吴，杜预随后接掌荆州军事部署}。\footnote{核心来源：《晋书·羊祜传》；《晋书·杜预传》；《资治通鉴》晋纪。材料限度：羊祜卒、表请伐吴和杜预继任大体可靠；表文传本、羊祜战略影响和朝廷决策先后有争议}

\paragraph{咸宁五年十一月（279年）} 沿着咸宁五年十一月（279年）留下的纪年节点，可以看到资源和权力怎样被逐项调度：\eventanchor{TK-0279-JIN-INVADES-WU}{西晋以杜预、王濬、王浑等分路伐吴，益州水军沿江东下，荆淮军并进}。\footnote{核心来源：《晋书·武帝纪》；《晋书·杜预传》；《晋书·王濬传》；《三国志·孙皓传》；《资治通鉴》晋纪。材料限度：出师时间、主要统帅和多路部署可靠；兵力数字、路线细节、各军协同与王濬舰队规模有争议}

\paragraph{太康元年三月（280年）} 同年的记载没有替行动者说明动机，却留下了可接续的顺序：\eventanchor{TK-0280-SUN-HAO-SURRENDER}{王濬舰队抵建业，孙皓出降，孙吴灭亡，西晋完成全国统一}。\footnote{核心来源：《晋书·武帝纪》；《晋书·王濬传》；《三国志·孙皓传》；《建康实录》。材料限度：建业降附、孙吴灭亡和统一结果可靠；王濬与王浑的功绩争议、降仪细节、战后安置和统一的实际覆盖范围有争议}


\subsection{晋纪所见：接收之前的制度与边地}

《晋书·武帝纪》的短条把受禅后的任命、赈恤、礼制、边地冲突和伐吴准备逐项留在年序里。它们显示接收仍在进行，不证明各郡或降众已以同一速度进入西晋秩序。

\paragraph{265年} 《晋书·武帝纪》在这一年并列留下赋役、赈济与赏赐的记录。\footnote{核心来源：《晋书·武帝纪》。材料限度：本纪可证实所载命令、任免或结果；不据此推定各郡执行、全部参与者或私人动机。；本纪选择朝廷可见事项；不得由一条记录推定各郡社会状况或具体群体动机。；本纪所载多为朝廷命令或结果，不足以证明地方执行、总量和受影响户数。；本纪以本政权视角记录对手或边缘群体，不足以断定对方内部决策、疆界或长期归属。}
\noindent\eventanchor{TK-SRC-0265-01}{《晋书·武帝纪》在265年记下“」按：上文既言「皇弟」，則作「機」爲是，通鑑七九正作「機」，今據改”。晋廷有司面对的并非抽象的朝局，而是洛阳的晋廷及其州郡网络这处可调度的节点；名号、岗位或仪式在洛阳的晋廷及其州郡网络被重新接续，随后谁能据此调动人手仍要看后续记录。应把它限定为本纪可证实所载命令、任免或结果；不据此推定各郡执行、全部参与者或私人动机。}\par\smallskip
\noindent\eventanchor{TK-SRC-0265-02}{从《晋书·武帝纪》的265年条目可见：“」帝初以禮讓，魏朝公卿何曾、王沈等固請，乃從之”。这句话把使者与受命者放在洛阳的晋廷及其州郡网络，使处理此事的空间条件不再被抹平；文书或使节抵达后，洛阳的晋廷及其州郡网络与外部的关系获得了一个新的、仍待后续行动检验的接口。在证据上能站稳的说法是本纪可证实所载命令、任免或结果；不据此推定各郡执行、全部参与者或私人动机。}\par\smallskip
\noindent\eventanchor{TK-SRC-0265-03}{“』炎虔奉皇運，寅畏天威，敬簡元辰，升壇受禪，告類上帝，永答眾望”是《晋书·武帝纪》留给265年的可核文字。它指向洛阳的晋廷及其州郡网络中的西晋纪事中的当事人，而不是一则可以脱离地点转述的消息；因此后续谈判、质问或防务安排有了可追溯的起点；关系的深浅却不能由一条记录替代。这则本纪记录所能保证的是本纪可证实所载命令、任免或结果；不据此推定各郡执行、全部参与者或私人动机。}\par\smallskip
\noindent\eventanchor{TK-SRC-0265-04}{《晋书·武帝纪》把265年的这一刻写作“按宣五王傳，清惠亭侯京薨，以文帝子機爲嗣，封燕王”。沿着这条记录，晋廷有司在洛阳的晋廷及其州郡网络的处置才有可查的落点；此人的退出改变了洛阳的晋廷及其州郡网络能够继续发令或守备的入口。这不是对全局的透视，而是本纪选择朝廷可见事项；不得由一条记录推定各郡社会状况或具体群体动机。}\par\smallskip
\noindent\eventanchor{TK-SRC-0265-05}{在《晋书·武帝纪》所编的年序里，265年出现“八月辛卯，文帝崩，太子嗣相國、晉王位”。西晋纪事中的当事人由此进入洛阳的晋廷及其州郡网络的具体事务，不能只被写成政权名下的符号；死亡或处置随即把洛阳的晋廷及其州郡网络的名位与所属人手重新摆上日程。这里不宜越过的界线是本纪选择朝廷可见事项；不得由一条记录推定各郡社会状况或具体群体动机。}\par\smallskip
\noindent\eventanchor{TK-SRC-0265-06}{查《晋书·武帝纪》，265年所载原文为“除舊嫌，解禁錮，亡官失爵者悉復之”。原文所见的行动者是晋廷有司，其所处的接口是洛阳的晋廷及其州郡网络；它使洛阳的晋廷及其州郡网络的名位或物资流向出现可查的调整，却不能折算为全境的户数与总量。这条文字可靠地指出本纪可证实所载命令、任免或结果；不据此推定各郡执行、全部参与者或私人动机。}\par\smallskip
\noindent\eventanchor{TK-SRC-0265-07}{265年的《晋书·武帝纪》没有略过这件事：“此舉其字者，蓋沛王名本或作『炳』及『昞』、『秉』之類，唐避世祖諱昞，於『丙』、『秉』等字皆改爲『景』”。它把晋廷有司与洛阳的晋廷及其州郡网络并置，留下了行动发生在哪里的证据边界；它让洛阳的晋廷及其州郡网络的权责关系出现新的书面形状，真正的落实范围不能从这一句外推。文本为我们保留的是本纪可证实所载命令、任免或结果；不据此推定各郡执行、全部参与者或私人动机。}\par\smallskip
\noindent\eventanchor{TK-SRC-0265-08}{《晋书·武帝纪》在265年保存“賜山陽公劉康、安樂公劉禪子弟一人爲駙馬都尉”一语；由此能够把劉禪的处置系在山陽，而非系在事后的概括上；其直接后果是山陽的赏赐、征收或救助次序被重新表达；执行规模仍留在材料之外。史料足以固定的范围是本纪选择朝廷可见事项；不得由一条记录推定各郡社会状况或具体群体动机。}\par\smallskip
\noindent\eventanchor{TK-SRC-0265-09}{“賜天下爵，人五級；鰥寡孤獨不能自存者穀，人五斛”见于《晋书·武帝纪》的265年记录。对晋廷有司而言，洛阳的晋廷及其州郡网络不是背景：它决定了这项动作要经由哪一层机构、道路或人群；资源经由洛阳的晋廷及其州郡网络的这一分配被标出优先次序，受及何人、实物多少仍须另有账簿才能落实。这项记述至多能够说明本纪所载多为朝廷命令或结果，不足以证明地方执行、总量和受影响户数。}\par\smallskip
\noindent\eventanchor{TK-SRC-0265-10}{《晋书·武帝纪》为265年留下的句子是“誕惟四方，罔不祗順，廓清梁岷，包懷揚越，八紘同軌，祥瑞屢臻，天人協應，無思不服”。据此可见西晋纪事中的当事人在洛阳的晋廷及其州郡网络面临的实际接口，而不必替其补写材料未载的盘算；这项制度性处置重排了洛阳的晋廷及其州郡网络可见的名位与入口，却不自动证明各地已按同一尺度执行。材料在此允许我们说的只是本纪可证实所载命令、任免或结果；不据此推定各郡执行、全部参与者或私人动机。}\par\smallskip
\noindent\eventanchor{TK-SRC-0265-11}{把目光放回《晋书·武帝纪》，265年的文字写作“丁卯，遣太僕劉原告于太廟”。使者与受命者在洛阳的晋廷及其州郡网络的这一动作因而有了可复核的出处、地点和时序；因此后续谈判、质问或防务安排有了可追溯的起点；关系的深浅却不能由一条记录替代。它首先证明的是本纪可证实所载命令、任免或结果；不据此推定各郡执行、全部参与者或私人动机。}\par\smallskip
\noindent\eventanchor{TK-SRC-0265-12}{265年，《晋书·武帝纪》以“方軌虞夏四代之明顯，我不敢知”收录此事。它让西晋纪事中的当事人从洛阳的晋廷及其州郡网络的制度、军路或文书网络中显形，而非漂浮在宏大叙事之外；这使洛阳的晋廷及其州郡网络的一段关系或制度调整有了可追踪的起点，代价与范围仍不可臆断。可据此确认的是本纪可证实所载命令、任免或结果；不据此推定各郡执行、全部参与者或私人动机。}\par\smallskip
\noindent\eventanchor{TK-SRC-0265-13}{就265年而言，《晋书·武帝纪》所能明确的是“封皇叔祖父孚爲安平王，皇叔父幹爲平原王，亮爲扶風王，伷爲東莞王，駿爲汝陰王，肜爲梁王，倫爲琅邪王，皇弟攸爲齊王，鑒爲樂安王，機爲燕王，機爲燕王「機」，各本皆作「幾」”。这把晋廷有司与洛阳的晋廷及其州郡网络连在同一条可检的记录上；资源经由洛阳的晋廷及其州郡网络的这一分配被标出优先次序，受及何人、实物多少仍须另有账簿才能落实。可从中取得的最低结论是本纪以本政权视角记录对手或边缘群体，不足以断定对方内部决策、疆界或长期归属。}\par\smallskip
\noindent\eventanchor{TK-SRC-0265-14}{《晋书·武帝纪》的265年纪事写下“封魏帝爲陳留王，邑萬戶，居於鄴宮；魏氏諸王皆爲縣侯”；文字所触及的是鄴中的晋廷有司，不是一个可以任意延展的全境结论；它使鄴的名位或物资流向出现可查的调整，却不能折算为全境的户数与总量。可核验的底线仍是本纪选择朝廷可见事项；不得由一条记录推定各郡社会状况或具体群体动机。}\par\smallskip
\noindent\eventanchor{TK-SRC-0265-15}{读《晋书·武帝纪》的265年条目，首先遇到的是“復天下租賦及關市之稅一年，逋債宿負皆勿收”。它说明西晋纪事中的当事人必须在洛阳的晋廷及其州郡网络的条件下作出这一处置，后果也从那里展开；这笔给付或减免给洛阳的晋廷及其州郡网络的资源组织增添了一个节点，不可据此假定同样措施已普遍兑现。读到这里须收住判断：本纪所载多为朝廷命令或结果，不足以证明地方执行、总量和受影响户数。}\par\smallskip
\noindent\eventanchor{TK-SRC-0265-16}{265年这一条由《晋书·武帝纪》记为“改景初曆爲太始曆，臘以酉，社以丑”。晋廷有司所处的洛阳的晋廷及其州郡网络由此成为理解动作与代价的起点；因此洛阳的晋廷及其州郡网络的下一步任命、礼仪或行政衔接获得了条件，而不是已经完成的结果。因此应严格区分：本纪可证实所载命令、任免或结果；不据此推定各郡执行、全部参与者或私人动机。}\par\smallskip
\noindent\eventanchor{TK-SRC-0265-17}{《晋书·武帝纪》在265年记下“格爾上下神祇，罔不克順，地平天成，萬邦以乂”。西晋纪事中的当事人面对的并非抽象的朝局，而是洛阳的晋廷及其州郡网络这处可调度的节点；这项处置为洛阳的晋廷及其州郡网络的后续任命、征发或回应留下了一个具体的时间节点。应把它限定为本纪可证实所载命令、任免或结果；不据此推定各郡执行、全部参与者或私人动机。}\par\smallskip
\noindent\eventanchor{TK-SRC-0265-18}{从《晋书·武帝纪》的265年条目可见：“故紀稱其字，傳改爲『景』耳”。这句话把晋廷有司放在洛阳的晋廷及其州郡网络，使处理此事的空间条件不再被抹平；这项制度性处置重排了洛阳的晋廷及其州郡网络可见的名位与入口，却不自动证明各地已按同一尺度执行。在证据上能站稳的说法是本纪可证实所载命令、任免或结果；不据此推定各郡执行、全部参与者或私人动机。}\par\smallskip
\noindent\eventanchor{TK-SRC-0265-19}{“薨，無子，齊王冏表以子幾嗣”是《晋书·武帝纪》留给265年的可核文字。它指向洛阳的晋廷及其州郡网络中的西晋纪事中的当事人，而不是一则可以脱离地点转述的消息；它留下的直接压力，是由谁接住洛阳的晋廷及其州郡网络中断的名义与调度。这则本纪记录所能保证的是本纪选择朝廷可见事项；不得由一条记录推定各郡社会状况或具体群体动机。}\par\smallskip
\noindent\eventanchor{TK-SRC-0265-20}{《晋书·武帝纪》把265年的这一刻写作“皇從伯父望爲義陽王，皇從叔父輔爲渤海王，晃爲下邳王，瑰爲太原王，珪爲高陽王，衡爲常山王，子文爲沛王，子文爲沛王李校：「宗室傳言沛王景字子文”。沿着这条记录，西晋纪事中的当事人在洛阳的晋廷及其州郡网络的处置才有可查的落点；这使洛阳的晋廷及其州郡网络的一段关系或制度调整有了可追踪的起点，代价与范围仍不可臆断。这不是对全局的透视，而是本纪选择朝廷可见事项；不得由一条记录推定各郡社会状况或具体群体动机。}\par\smallskip
\noindent\eventanchor{TK-SRC-0265-21}{在《晋书·武帝纪》所编的年序里，265年出现“己巳，詔陳留王載天子旌旗，備五時副車，行魏正朔，郊祀天地，禮樂制度皆如魏舊，上書不稱臣”。晋廷有司由此进入洛阳的晋廷及其州郡网络的具体事务，不能只被写成政权名下的符号；这次往来把洛阳的晋廷及其州郡网络纳入一条可以继续追踪的联系线，却不等于长期服从已经形成。这里不宜越过的界线是本纪选择朝廷可见事项；不得由一条记录推定各郡社会状况或具体群体动机。}\par\smallskip
\noindent\eventanchor{TK-SRC-0265-22}{查《晋书·武帝纪》，265年所载原文为“暨漢德既衰，太祖武皇帝撥亂濟時，扶翼劉氏，又用受命于漢”。原文所见的行动者是西晋纪事中的当事人，其所处的接口是洛阳的晋廷及其州郡网络；文书或使节抵达后，洛阳的晋廷及其州郡网络与外部的关系获得了一个新的、仍待后续行动检验的接口。这条文字可靠地指出本纪以本政权视角记录对手或边缘群体，不足以断定对方内部决策、疆界或长期归属。}\par\smallskip
\noindent\eventanchor{TK-SRC-0265-23}{265年的《晋书·武帝纪》没有略过这件事：“九月戊午，以魏司徒何曾爲丞相，鎮南將軍王沈爲御史大夫，中護軍賈充爲衛將軍，議郎裴秀爲尚書令、光祿大夫，皆開府”。它把賈充与洛阳的晋廷及其州郡网络并置，留下了行动发生在哪里的证据边界；这项行动把后果落在洛阳的晋廷及其州郡网络的通行与防守上，而非只改变一纸战报。文本为我们保留的是本纪以本政权视角记录对手或边缘群体，不足以断定对方内部决策、疆界或长期归属。}\par\smallskip
\noindent\eventanchor{TK-SRC-0265-24}{《晋书·武帝纪》在265年保存“率循訓典，底綏四國，用保天休，無替我二皇之弘烈”一语；由此能够把西晋纪事中的当事人的处置系在洛阳的晋廷及其州郡网络，而非系在事后的概括上；这使洛阳的晋廷及其州郡网络的一段关系或制度调整有了可追踪的起点，代价与范围仍不可臆断。史料足以固定的范围是本纪可证实所载命令、任免或结果；不据此推定各郡执行、全部参与者或私人动机。}\par\smallskip
\noindent\eventanchor{TK-SRC-0265-25}{“其餘增封進爵各有差，文武普增位二等”见于《晋书·武帝纪》的265年记录。对晋廷有司而言，洛阳的晋廷及其州郡网络不是背景：它决定了这项动作要经由哪一层机构、道路或人群；资源经由洛阳的晋廷及其州郡网络的这一分配被标出优先次序，受及何人、实物多少仍须另有账簿才能落实。这项记述至多能够说明本纪以本政权视角记录对手或边缘群体，不足以断定对方内部决策、疆界或长期归属。}\par\smallskip
\noindent\eventanchor{TK-SRC-0265-26}{《晋书·武帝纪》为265年留下的句子是“十一月，初置四護軍，以統城外諸軍”。据此可见晋廷有司在洛阳的晋廷及其州郡网络面临的实际接口，而不必替其补写材料未载的盘算；军队一旦进入这一节点，洛阳的晋廷及其州郡网络的道路、水面或城防便承受了后续调度的压力。材料在此允许我们说的只是本纪以本政权视角记录对手或边缘群体，不足以断定对方内部决策、疆界或长期归属。}\par\smallskip
\noindent\eventanchor{TK-SRC-0265-27}{把目光放回《晋书·武帝纪》，265年的文字写作“是月，鳳皇六、青龍三、白龍二、麒麟各一見于郡國”。西晋纪事中的当事人在洛阳的晋廷及其州郡网络的这一动作因而有了可复核的出处、地点和时序；由此留下的是洛阳的晋廷及其州郡网络在文本中的异常时点，而非可量化的社会后果。它首先证明的是本纪可证实所载命令、任免或结果；不据此推定各郡执行、全部参与者或私人动机。}\par\smallskip
\noindent\eventanchor{TK-SRC-0265-28}{265年，《晋书·武帝纪》以“肆予憲章三后，用集大命于茲”收录此事。它让西晋纪事中的当事人从洛阳的晋廷及其州郡网络的制度、军路或文书网络中显形，而非漂浮在宏大叙事之外；这使洛阳的晋廷及其州郡网络的一段关系或制度调整有了可追踪的起点，代价与范围仍不可臆断。可据此确认的是本纪选择朝廷可见事项；不得由一条记录推定各郡社会状况或具体群体动机。}\par\smallskip
\noindent\eventanchor{TK-SRC-0265-29}{就265年而言，《晋书·武帝纪》所能明确的是“肆予一人，祗承天序，以敬授爾位，曆數實在爾躬”。这把西晋纪事中的当事人与洛阳的晋廷及其州郡网络连在同一条可检的记录上；这项处置为洛阳的晋廷及其州郡网络的后续任命、征发或回应留下了一个具体的时间节点。可从中取得的最低结论是本纪可证实所载命令、任免或结果；不据此推定各郡执行、全部参与者或私人动机。}\par\smallskip
\noindent\eventanchor{TK-SRC-0265-30}{《晋书·武帝纪》的265年纪事写下“泰爲隴西王，權爲彭城王，綏爲范陽王，遂爲濟南王，遜爲譙王，睦爲中山王，陵爲北海王，斌爲陳王，皇從父兄洪爲河間王，皇從父弟楙爲東平王”；文字所触及的是洛阳的晋廷及其州郡网络中的西晋纪事中的当事人，不是一个可以任意延展的全境结论；它把洛阳的晋廷及其州郡网络既有的安排推向下一步，结果的广度仍须由别的材料补证。可核验的底线仍是本纪以本政权视角记录对手或边缘群体，不足以断定对方内部决策、疆界或长期归属。}\par\smallskip
\noindent\eventanchor{TK-SRC-0265-31}{读《晋书·武帝纪》的265年条目，首先遇到的是“天序不可以無統，人神不可以曠主”。它说明西晋纪事中的当事人必须在洛阳的晋廷及其州郡网络的条件下作出这一处置，后果也从那里展开；由此可接续观察洛阳的晋廷及其州郡网络的后续变化，但不能把单项记录扩充为完整的政治图景。读到这里须收住判断：本纪可证实所载命令、任免或结果；不据此推定各郡执行、全部参与者或私人动机。}\par\smallskip
\noindent\eventanchor{TK-SRC-0265-32}{265年这一条由《晋书·武帝纪》记为“惟三后陟配於天，而咸用光敷聖德”。西晋纪事中的当事人所处的洛阳的晋廷及其州郡网络由此成为理解动作与代价的起点；这使洛阳的晋廷及其州郡网络的一段关系或制度调整有了可追踪的起点，代价与范围仍不可臆断。因此应严格区分：本纪选择朝廷可见事项；不得由一条记录推定各郡社会状况或具体群体动机。}\par\smallskip
\noindent\eventanchor{TK-SRC-0265-33}{《晋书·武帝纪》在265年记下“惟王乃祖乃父，服膺明哲，輔亮我皇家，勳德光於四海”。西晋纪事中的当事人面对的并非抽象的朝局，而是洛阳的晋廷及其州郡网络这处可调度的节点；名号、岗位或仪式在洛阳的晋廷及其州郡网络被重新接续，随后谁能据此调动人手仍要看后续记录。应把它限定为本纪以本政权视角记录对手或边缘群体，不足以断定对方内部决策、疆界或长期归属。}\par\smallskip
\noindent\eventanchor{TK-SRC-0265-34}{从《晋书·武帝纪》的265年条目可见：“下令寬刑宥罪，撫眾息役，國內行服三日”。这句话把西晋纪事中的当事人放在洛阳的晋廷及其州郡网络，使处理此事的空间条件不再被抹平；文书或使节抵达后，洛阳的晋廷及其州郡网络与外部的关系获得了一个新的、仍待后续行动检验的接口。在证据上能站稳的说法是本纪选择朝廷可见事项；不得由一条记录推定各郡社会状况或具体群体动机。}\par\smallskip
\noindent\eventanchor{TK-SRC-0265-35}{“咸熙二年五月，立爲晉王太子”是《晋书·武帝纪》留给265年的可核文字。它指向洛阳的晋廷及其州郡网络中的西晋纪事中的当事人，而不是一则可以脱离地点转述的消息；它让洛阳的晋廷及其州郡网络的权责关系出现新的书面形状，真正的落实范围不能从这一句外推。这则本纪记录所能保证的是本纪选择朝廷可见事项；不得由一条记录推定各郡社会状况或具体群体动机。}\par\smallskip
\noindent\eventanchor{TK-SRC-0265-36}{《晋书·武帝纪》把265年的这一刻写作“乙亥，以安平王孚爲太宰、假黃鉞、大都督中外諸軍事”。沿着这条记录，西晋的将领与所部军队在洛阳的晋廷及其州郡网络的处置才有可查的落点；它迫使相关一方在洛阳的晋廷及其州郡网络重新估量驻屯、救援和退路的次序。这不是对全局的透视，而是本纪以本政权视角记录对手或边缘群体，不足以断定对方内部决策、疆界或长期归属。}\par\smallskip
\noindent\eventanchor{TK-SRC-0265-37}{在《晋书·武帝纪》所编的年序里，265年出现“乙未，令諸郡中正以六條舉淹滯：一曰忠恪匪躬，二曰孝敬盡禮，三曰友于兄弟，四曰潔身勞謙，五曰信義可復，六曰學以爲己”。西晋纪事中的当事人由此进入洛阳的晋廷及其州郡网络的具体事务，不能只被写成政权名下的符号；这次往来把洛阳的晋廷及其州郡网络纳入一条可以继续追踪的联系线，却不等于长期服从已经形成。这里不宜越过的界线是本纪选择朝廷可见事项；不得由一条记录推定各郡社会状况或具体群体动机。}\par\smallskip
\noindent\eventanchor{TK-SRC-0265-38}{查《晋书·武帝纪》，265年所载原文为“以驃騎將軍石苞爲大司馬，封樂陵公，車騎將軍陳騫爲高平公，衛將軍賈充爲車騎將軍、魯公，尚書令裴秀爲鉅鹿公，侍中荀勖爲濟北公，太保鄭沖爲太傅、壽光公，太尉王祥爲太保、睢陵公，丞相何曾爲太尉、朗陵公，御史大夫王沈爲驃騎將軍、博陵公，司空荀顗爲臨淮公，鎮北大將軍衛瓘爲菑陽公”。原文所见的行动者是賈充，其所处的接口是洛阳的晋廷及其州郡网络；军队一旦进入这一节点，洛阳的晋廷及其州郡网络的道路、水面或城防便承受了后续调度的压力。这条文字可靠地指出本纪以本政权视角记录对手或边缘群体，不足以断定对方内部决策、疆界或长期归属。}\par\smallskip
\noindent\eventanchor{TK-SRC-0265-39}{265年的《晋书·武帝纪》没有略过这件事：“應受上帝之命，協皇極之中”。它把西晋纪事中的当事人与洛阳的晋廷及其州郡网络并置，留下了行动发生在哪里的证据边界；因此后续谈判、质问或防务安排有了可追溯的起点；关系的深浅却不能由一条记录替代。文本为我们保留的是本纪以本政权视角记录对手或边缘群体，不足以断定对方内部决策、疆界或长期归属。}\par\smallskip
\noindent\eventanchor{TK-SRC-0265-40}{《晋书·武帝纪》在265年保存“置中軍將軍，以統宿衛七軍”一语；由此能够把晋廷有司的处置系在洛阳的晋廷及其州郡网络，而非系在事后的概括上；它迫使相关一方在洛阳的晋廷及其州郡网络重新估量驻屯、救援和退路的次序。史料足以固定的范围是本纪以本政权视角记录对手或边缘群体，不足以断定对方内部决策、疆界或长期归属。}\par\smallskip
\noindent\eventanchor{TK-SRC-0265-41}{“追尊宣王爲宣皇帝，景王爲景皇帝，文王爲文皇帝，宣王妃張氏爲宣穆皇后”见于《晋书·武帝纪》的265年记录。对西晋纪事中的当事人而言，洛阳的晋廷及其州郡网络不是背景：它决定了这项动作要经由哪一层机构、道路或人群；这项处置为洛阳的晋廷及其州郡网络的后续任命、征发或回应留下了一个具体的时间节点。这项记述至多能够说明本纪选择朝廷可见事项；不得由一条记录推定各郡社会状况或具体群体动机。}\par\smallskip
\noindent\eventanchor{TK-SRC-0265-42}{《晋书·武帝纪》为265年留下的句子是“尊太妃王氏曰皇太后，宮曰崇化”。据此可见西晋纪事中的当事人在洛阳的晋廷及其州郡网络面临的实际接口，而不必替其补写材料未载的盘算；这项制度性处置重排了洛阳的晋廷及其州郡网络可见的名位与入口，却不自动证明各地已按同一尺度执行。材料在此允许我们说的只是本纪选择朝廷可见事项；不得由一条记录推定各郡社会状况或具体群体动机。}\par\smallskip

\paragraph{266年} 《晋书·武帝纪》在这一年并列留下官署、任命与诏令的记录。\footnote{核心来源：《晋书·武帝纪》。材料限度：本纪可核对行动的基本顺序；兵额、路线、战损和具体指挥过程仍须另证。；本纪以本政权视角记录对手或边缘群体，不足以断定对方内部决策、疆界或长期归属。；本纪选择朝廷可见事项；不得由一条记录推定各郡社会状况或具体群体动机。}
\noindent\eventanchor{TK-SRC-0266-01}{把目光放回《晋书·武帝纪》，266年的文字写作“罷山陽公國督軍，除其禁制”。晋廷有司在山陽的这一动作因而有了可复核的出处、地点和时序；这项行动把后果落在山陽的通行与防守上，而非只改变一纸战报。它首先证明的是本纪可核对行动的基本顺序；兵额、路线、战损和具体指挥过程仍须另证。}\par\smallskip
\noindent\eventanchor{TK-SRC-0266-02}{266年，《晋书·武帝纪》以“并圜丘、方丘於南、北郊，二至之祀合於二郊”收录此事。它让西晋纪事中的当事人从洛阳的晋廷及其州郡网络的制度、军路或文书网络中显形，而非漂浮在宏大叙事之外；因此洛阳的晋廷及其州郡网络的下一步任命、礼仪或行政衔接获得了条件，而不是已经完成的结果。可据此确认的是本纪以本政权视角记录对手或边缘群体，不足以断定对方内部决策、疆界或长期归属。}\par\smallskip
\noindent\eventanchor{TK-SRC-0266-03}{就266年而言，《晋书·武帝纪》所能明确的是“初，帝雖從漢魏之制，既葬除服，而深衣素冠，降席撤膳，哀敬如喪者”。这把晋廷有司与洛阳的晋廷及其州郡网络连在同一条可检的记录上；这次往来把洛阳的晋廷及其州郡网络纳入一条可以继续追踪的联系线，却不等于长期服从已经形成。可从中取得的最低结论是本纪选择朝廷可见事项；不得由一条记录推定各郡社会状况或具体群体动机。}\par\smallskip
\noindent\eventanchor{TK-SRC-0266-04}{《晋书·武帝纪》的266年纪事写下“丁亥，有司請建七廟，帝重其役，不許”；文字所触及的是洛阳的晋廷及其州郡网络中的晋廷有司，不是一个可以任意延展的全境结论；文书或使节抵达后，洛阳的晋廷及其州郡网络与外部的关系获得了一个新的、仍待后续行动检验的接口。可核验的底线仍是本纪选择朝廷可见事项；不得由一条记录推定各郡社会状况或具体群体动机。}\par\smallskip
\noindent\eventanchor{TK-SRC-0266-05}{读《晋书·武帝纪》的266年条目，首先遇到的是“冬十月，鮮卑慕容廆寇昌黎”。它说明慕容廆必须在昌黎的条件下作出这一处置，后果也从那里展开；这项行动把后果落在昌黎的通行与防守上，而非只改变一纸战报。读到这里须收住判断：本纪以本政权视角记录对手或边缘群体，不足以断定对方内部决策、疆界或长期归属。}\par\smallskip
\noindent\eventanchor{TK-SRC-0266-06}{266年这一条由《晋书·武帝纪》记为“二年春二月，淮南、丹楊地震”。当地官吏与受灾住民所处的淮南由此成为理解动作与代价的起点；它改变的是淮南被书写和被解释的方式，实际影响仍不应由灾异文字代替。因此应严格区分：本纪以本政权视角记录对手或边缘群体，不足以断定对方内部决策、疆界或长期归属。}\par\smallskip
\noindent\eventanchor{TK-SRC-0266-07}{《晋书·武帝纪》在266年记下“二年春正月丙戌，遣兼侍中侯史光等持節四方，循省風俗，除禳祝之不在祀典者”。晋廷有司面对的并非抽象的朝局，而是洛阳的晋廷及其州郡网络这处可调度的节点；这次往来把洛阳的晋廷及其州郡网络纳入一条可以继续追踪的联系线，却不等于长期服从已经形成。应把它限定为本纪以本政权视角记录对手或边缘群体，不足以断定对方内部决策、疆界或长期归属。}\par\smallskip
\noindent\eventanchor{TK-SRC-0266-08}{从《晋书·武帝纪》的266年条目可见：“江夏、泰山水，流居人三百餘家”。这句话把当地官吏与受灾住民放在江夏，使处理此事的空间条件不再被抹平；军队一旦进入这一节点，江夏的道路、水面或城防便承受了后续调度的压力。在证据上能站稳的说法是本纪以本政权视角记录对手或边缘群体，不足以断定对方内部决策、疆界或长期归属。}\par\smallskip

\paragraph{267年} 《晋书·武帝纪》在这一年并列留下赋役、赈济与赏赐的记录。\footnote{核心来源：《晋书·武帝纪》。材料限度：本纪所载多为朝廷命令或结果，不足以证明地方执行、总量和受影响户数。；本纪可核对行动的基本顺序；兵额、路线、战损和具体指挥过程仍须另证。}
\noindent\eventanchor{TK-SRC-0267-01}{“丙申，詔四方水旱甚者無出田租”是《晋书·武帝纪》留给267年的可核文字。它指向洛阳的晋廷及其州郡网络中的晋廷有司，而不是一则可以脱离地点转述的消息；因此后续谈判、质问或防务安排有了可追溯的起点；关系的深浅却不能由一条记录替代。这则本纪记录所能保证的是本纪所载多为朝廷命令或结果，不足以证明地方执行、总量和受影响户数。}\par\smallskip
\noindent\eventanchor{TK-SRC-0267-02}{《晋书·武帝纪》把267年的这一刻写作“秋八月，罷都護將軍，以其五署還光祿勳”。沿着这条记录，晋廷有司在洛阳的晋廷及其州郡网络的处置才有可查的落点；它迫使相关一方在洛阳的晋廷及其州郡网络重新估量驻屯、救援和退路的次序。这不是对全局的透视，而是本纪可核对行动的基本顺序；兵额、路线、战损和具体指挥过程仍须另证。}\par\smallskip
\noindent\eventanchor{TK-SRC-0267-03}{在《晋书·武帝纪》所编的年序里，267年出现“吳故將莞恭、帛奉舉兵反，攻害建鄴令，遂圍揚州，徐州刺史嵇喜討平之”。嵇喜由此进入徐州的具体事务，不能只被写成政权名下的符号；由此，徐州的守备、转输与进退选择须随眼前的军情重新安排。这里不宜越过的界线是本纪所载多为朝廷命令或结果，不足以证明地方执行、总量和受影响户数。}\par\smallskip
\noindent\eventanchor{TK-SRC-0267-04}{查《晋书·武帝纪》，267年所载原文为“夏四月戊午，張掖太守焦勝上言，氐池縣大柳谷口有玄石一所，白畫成文，實大晉之休祥，圖之以獻”。原文所见的行动者是焦勝，其所处的接口是張掖；文书或使节抵达后，張掖与外部的关系获得了一个新的、仍待后续行动检验的接口。这条文字可靠地指出本纪可核对行动的基本顺序；兵额、路线、战损和具体指挥过程仍须另证。}\par\smallskip

\paragraph{268年} 《晋书·武帝纪》在这一年并列留下赋役、赈济与赏赐的记录。\footnote{核心来源：《晋书·武帝纪》。材料限度：本纪可核对行动的基本顺序；兵额、路线、战损和具体指挥过程仍须另证。；本纪所载多为朝廷命令或结果，不足以证明地方执行、总量和受影响户数。}
\noindent\eventanchor{TK-SRC-0268-01}{268年的《晋书·武帝纪》没有略过这件事：“罷振威、揚威護軍官，置左右積弩將軍”。它把晋廷有司与洛阳的晋廷及其州郡网络并置，留下了行动发生在哪里的证据边界；这项行动把后果落在洛阳的晋廷及其州郡网络的通行与防守上，而非只改变一纸战报。文本为我们保留的是本纪可核对行动的基本顺序；兵额、路线、战损和具体指挥过程仍须另证。}\par\smallskip
\noindent\eventanchor{TK-SRC-0268-02}{《晋书·武帝纪》在268年保存“本書食貨志及御覽一九〇引晉陽秋，謂泰始四年七月起常平倉”一语；由此能够把西晋纪事中的当事人的处置系在洛阳的晋廷及其州郡网络，而非系在事后的概括上；它改变的是洛阳的晋廷及其州郡网络上可被承认的接触方式，而非替任何一方预先写定忠顺与疆界。史料足以固定的范围是本纪所载多为朝廷命令或结果，不足以证明地方执行、总量和受影响户数。}\par\smallskip
\noindent\eventanchor{TK-SRC-0268-03}{“丙戌，律令成，封爵賜帛各有差”见于《晋书·武帝纪》的268年记录。对晋廷有司而言，洛阳的晋廷及其州郡网络不是背景：它决定了这项动作要经由哪一层机构、道路或人群；这次往来把洛阳的晋廷及其州郡网络纳入一条可以继续追踪的联系线，却不等于长期服从已经形成。这项记述至多能够说明本纪所载多为朝廷命令或结果，不足以证明地方执行、总量和受影响户数。}\par\smallskip
\noindent\eventanchor{TK-SRC-0268-04}{《晋书·武帝纪》为268年留下的句子是“丙寅，兗州大水，復其田租”。据此可见当地官吏与受灾住民在兗州面临的实际接口，而不必替其补写材料未载的盘算；它使兗州的名位或物资流向出现可查的调整，却不能折算为全境的户数与总量。材料在此允许我们说的只是本纪所载多为朝廷命令或结果，不足以证明地方执行、总量和受影响户数。}\par\smallskip
\noindent\eventanchor{TK-SRC-0268-05}{把目光放回《晋书·武帝纪》，268年的文字写作“己未，詔王公卿尹及郡國守相，舉賢良方正直言之士”。晋廷有司在洛阳的晋廷及其州郡网络的这一动作因而有了可复核的出处、地点和时序；因此后续谈判、质问或防务安排有了可追溯的起点；关系的深浅却不能由一条记录替代。它首先证明的是本纪可核对行动的基本顺序；兵额、路线、战损和具体指挥过程仍须另证。}\par\smallskip
\noindent\eventanchor{TK-SRC-0268-06}{268年，《晋书·武帝纪》以“九月，青、徐、兗、豫四州大水，伊洛溢，合於河，開倉以振之”收录此事。它让当地官吏与受灾住民从伊洛的制度、军路或文书网络中显形，而非漂浮在宏大叙事之外；其直接后果是伊洛的赏赐、征收或救助次序被重新表达；执行规模仍留在材料之外。可据此确认的是本纪所载多为朝廷命令或结果，不足以证明地方执行、总量和受影响户数。}\par\smallskip
\noindent\eventanchor{TK-SRC-0268-07}{就268年而言，《晋书·武帝纪》所能明确的是“是月，以平州刺史傅詢、前廣平太守孟桓清白有聞，詢賜帛二百匹，桓百匹”。这把傅詢与洛阳的晋廷及其州郡网络连在同一条可检的记录上；资源经由洛阳的晋廷及其州郡网络的这一分配被标出优先次序，受及何人、实物多少仍须另有账簿才能落实。可从中取得的最低结论是本纪所载多为朝廷命令或结果，不足以证明地方执行、总量和受影响户数。}\par\smallskip
\noindent\eventanchor{TK-SRC-0268-08}{《晋书·武帝纪》的268年纪事写下“五月己亥，大將軍、琅邪王伷薨”；文字所触及的是洛阳的晋廷及其州郡网络中的西晋的将领与所部军队，不是一个可以任意延展的全境结论；这一终止使洛阳的晋廷及其州郡网络原有的继承与军政衔接不能再被视为既成。可核验的底线仍是本纪可核对行动的基本顺序；兵额、路线、战损和具体指挥过程仍须另证。}\par\smallskip

\paragraph{269年} 《晋书·武帝纪》在这一年并列留下赋役、赈济与赏赐的记录。\footnote{核心来源：《晋书·武帝纪》。材料限度：本纪可核对行动的基本顺序；兵额、路线、战损和具体指挥过程仍须另证。；本纪所载多为朝廷命令或结果，不足以证明地方执行、总量和受影响户数。}
\noindent\eventanchor{TK-SRC-0269-01}{读《晋书·武帝纪》的269年条目，首先遇到的是“罷鎮軍將軍，復置左右將軍官”。它说明晋廷有司必须在洛阳的晋廷及其州郡网络的条件下作出这一处置，后果也从那里展开；这项行动把后果落在洛阳的晋廷及其州郡网络的通行与防守上，而非只改变一纸战报。读到这里须收住判断：本纪可核对行动的基本顺序；兵额、路线、战损和具体指挥过程仍须另证。}\par\smallskip
\noindent\eventanchor{TK-SRC-0269-02}{269年这一条由《晋书·武帝纪》记为“冬十月丙子，以汲郡太守王宏有政績，賜穀千斛”。晋廷有司所处的汲郡由此成为理解动作与代价的起点；其直接后果是汲郡的赏赐、征收或救助次序被重新表达；执行规模仍留在材料之外。因此应严格区分：本纪所载多为朝廷命令或结果，不足以证明地方执行、总量和受影响户数。}\par\smallskip
\noindent\eventanchor{TK-SRC-0269-03}{《晋书·武帝纪》在269年记下“二月，以雍州隴右五郡及涼州之金城、梁州之陰平置秦州”。晋廷有司面对的并非抽象的朝局，而是隴右这处可调度的节点；名号、岗位或仪式在隴右被重新接续，随后谁能据此调动人手仍要看后续记录。应把它限定为本纪可核对行动的基本顺序；兵额、路线、战损和具体指挥过程仍须另证。}\par\smallskip
\noindent\eventanchor{TK-SRC-0269-04}{从《晋书·武帝纪》的269年条目可见：“二月戊午，王濬、唐彬等克丹楊城”。这句话把王濬放在丹楊，使处理此事的空间条件不再被抹平；它把丹楊既有的安排推向下一步，结果的广度仍须由别的材料补证。在证据上能站稳的说法是本纪可核对行动的基本顺序；兵额、路线、战损和具体指挥过程仍须另证。}\par\smallskip
\noindent\eventanchor{TK-SRC-0269-05}{“庚辰，以王濬爲輔國大將軍、襄陽侯，杜預當陽侯，王戎安豐侯，唐彬上庸侯，賈充、琅邪王伷以下增封”是《晋书·武帝纪》留给269年的可核文字。它指向襄陽中的王濬、杜預，而不是一则可以脱离地点转述的消息；这项行动把后果落在襄陽的通行与防守上，而非只改变一纸战报。这则本纪记录所能保证的是本纪可核对行动的基本顺序；兵额、路线、战损和具体指挥过程仍须另证。}\par\smallskip
\noindent\eventanchor{TK-SRC-0269-06}{《晋书·武帝纪》把269年的这一刻写作“六月，鄴奚官督郭廙上疏陳五事以諫，言甚切直，擢爲屯留令”。沿着这条记录，郭廙在鄴的处置才有可查的落点；它迫使相关一方在鄴重新估量驻屯、救援和退路的次序。这不是对全局的透视，而是本纪可核对行动的基本顺序；兵额、路线、战损和具体指挥过程仍须另证。}\par\smallskip
\noindent\eventanchor{TK-SRC-0269-07}{在《晋书·武帝纪》所编的年序里，269年出现“六月丁丑，初置翊軍校尉官”。晋廷有司由此进入洛阳的晋廷及其州郡网络的具体事务，不能只被写成政权名下的符号；由此，洛阳的晋廷及其州郡网络的守备、转输与进退选择须随眼前的军情重新安排。这里不宜越过的界线是本纪可核对行动的基本顺序；兵额、路线、战损和具体指挥过程仍须另证。}\par\smallskip
\noindent\eventanchor{TK-SRC-0269-08}{查《晋书·武帝纪》，269年所载原文为“五年春正月癸巳，申戒郡國計吏守相令長，務盡地利，禁游食商販”。原文所见的行动者是西晋纪事中的当事人，其所处的接口是洛阳的晋廷及其州郡网络；文书或使节抵达后，洛阳的晋廷及其州郡网络与外部的关系获得了一个新的、仍待后续行动检验的接口。这条文字可靠地指出本纪所载多为朝廷命令或结果，不足以证明地方执行、总量和受影响户数。}\par\smallskip
\noindent\eventanchor{TK-SRC-0269-09}{269年的《晋书·武帝纪》没有略过这件事：“西平人麴路伐登聞鼓，言多祅謗，有司奏棄市”。它把晋廷有司与洛阳的晋廷及其州郡网络并置，留下了行动发生在哪里的证据边界；因此后续谈判、质问或防务安排有了可追溯的起点；关系的深浅却不能由一条记录替代。文本为我们保留的是本纪所载多为朝廷命令或结果，不足以证明地方执行、总量和受影响户数。}\par\smallskip
\noindent\eventanchor{TK-SRC-0269-10}{《晋书·武帝纪》在269年保存“夏四月，任城、魯國池水赤如血”一语；由此能够把当地官吏与受灾住民的处置系在任城，而非系在事后的概括上；因此任城的下一步任命、礼仪或行政衔接获得了条件，而不是已经完成的结果。史料足以固定的范围是本纪可核对行动的基本顺序；兵额、路线、战损和具体指挥过程仍须另证。}\par\smallskip
\noindent\eventanchor{TK-SRC-0269-11}{“以太尉賈充爲大都督，行冠軍將軍楊濟爲副，總統眾軍”见于《晋书·武帝纪》的269年记录。对賈充而言，洛阳的晋廷及其州郡网络不是背景：它决定了这项动作要经由哪一层机构、道路或人群；由此，洛阳的晋廷及其州郡网络的守备、转输与进退选择须随眼前的军情重新安排。这项记述至多能够说明本纪可核对行动的基本顺序；兵额、路线、战损和具体指挥过程仍须另证。}\par\smallskip
\noindent\eventanchor{TK-SRC-0269-12}{《晋书·武帝纪》为269年留下的句子是“於是論功行封，賜公卿以下帛各有差”。据此可见晋廷有司在洛阳的晋廷及其州郡网络面临的实际接口，而不必替其补写材料未载的盘算；它使洛阳的晋廷及其州郡网络的名位或物资流向出现可查的调整，却不能折算为全境的户数与总量。材料在此允许我们说的只是本纪所载多为朝廷命令或结果，不足以证明地方执行、总量和受影响户数。}\par\smallskip

\paragraph{270年} 《晋书·武帝纪》在这一年并列留下赋役、赈济与赏赐的记录。\footnote{核心来源：《晋书·武帝纪》。材料限度：本纪所载多为朝廷命令或结果，不足以证明地方执行、总量和受影响户数。；本纪可核对行动的基本顺序；兵额、路线、战损和具体指挥过程仍须另证。}
\noindent\eventanchor{TK-SRC-0270-02}{把目光放回《晋书·武帝纪》，270年的文字写作“冬十一月，幸辟雍，行鄉飲酒之禮，賜太常博士、學生帛牛酒各有差”。晋廷有司在洛阳的晋廷及其州郡网络的这一动作因而有了可复核的出处、地点和时序；这笔给付或减免给洛阳的晋廷及其州郡网络的资源组织增添了一个节点，不可据此假定同样措施已普遍兑现。它首先证明的是本纪所载多为朝廷命令或结果，不足以证明地方执行、总量和受影响户数。}\par\smallskip
\noindent\eventanchor{TK-SRC-0270-03}{270年，《晋书·武帝纪》以“秋七月丁酉，復隴右五郡遇寇害者租賦，不能自存者廩貸之”收录此事。它让西晋的将领与所部军队从隴右的制度、军路或文书网络中显形，而非漂浮在宏大叙事之外；它迫使相关一方在隴右重新估量驻屯、救援和退路的次序。可据此确认的是本纪所载多为朝廷命令或结果，不足以证明地方执行、总量和受影响户数。}\par\smallskip
\noindent\eventanchor{TK-SRC-0270-04}{就270年而言，《晋书·武帝纪》所能明确的是“以鎮軍大將軍王濬爲撫軍大將軍”。这把王濬与洛阳的晋廷及其州郡网络连在同一条可检的记录上；由此，洛阳的晋廷及其州郡网络的守备、转输与进退选择须随眼前的军情重新安排。可从中取得的最低结论是本纪可核对行动的基本顺序；兵额、路线、战损和具体指挥过程仍须另证。}\par\smallskip
\noindent\eventanchor{TK-SRC-0270-05}{《晋书·武帝纪》的270年纪事写下“詔遣尚書石鑒行安西將軍、都督秦州諸軍事，與奮威護軍田章討之”；文字所触及的是秦州中的晋廷有司，不是一个可以任意延展的全境结论；军队一旦进入这一节点，秦州的道路、水面或城防便承受了后续调度的压力。可核验的底线仍是本纪所载多为朝廷命令或结果，不足以证明地方执行、总量和受影响户数。}\par\smallskip

\paragraph{271年} 《晋书·武帝纪》在这一年并列留下官署、任命与诏令的记录。\footnote{核心来源：《晋书·武帝纪》。材料限度：本纪可核对行动的基本顺序；兵额、路线、战损和具体指挥过程仍须另证。}
\noindent\eventanchor{TK-SRC-0271-01}{读《晋书·武帝纪》的271年条目，首先遇到的是“六月，詔公卿以下舉將帥各一人”。它说明晋廷有司必须在洛阳的晋廷及其州郡网络的条件下作出这一处置，后果也从那里展开；因此后续谈判、质问或防务安排有了可追溯的起点；关系的深浅却不能由一条记录替代。读到这里须收住判断：本纪可核对行动的基本顺序；兵额、路线、战损和具体指挥过程仍须另证。}\par\smallskip
